\section{LlamaIndex 源码分析：Indexing 流程}

\subsection{文档加载与 Chunking}

\begin{enumerate}
    \item 从 Data 目录加载所有文档（返回 Document 列表）
    \item 对每个文档进行 \textbf{Split}——Chunk Size 为 1024（实际约 997，减去文档路径的 token 数），Overlap 为 200
    \item Split 过程先按自身规则切分（产生约 759 个 split），再进行 \textbf{Merge} 合并为最终的 Chunk（约 22 个）
\end{enumerate}

\subsection{Node 创建与 Embedding}

\begin{enumerate}
    \item 每个 Chunk 创建一个 \textbf{Node}（具有唯一 node\_id 和对应的文本内容）
    \item 使用 Embedding Model（text-embedding-ada-002，Batch Size 2048）为每个 Node 计算 Embedding Vector
    \item 最终得到 22 个 Node，每个 Node 有一个 1536 维的 Embedding Vector（$512 \times 3$）
\end{enumerate}

\subsection{本章小结}

Indexing 流程：文档 $\rightarrow$ Split $\rightarrow$ Merge $\rightarrow$ Chunk $\rightarrow$ Node $\rightarrow$ Embedding Vector。每一步都在源码中有清晰的实现。

